\section{Controlled Study Design}
\label{sec:method}
The study treats feature geometry and network capacity as separate experimental
factors. The corpus, observation length, class interface, data partitions,
optimizer family, learning-rate schedule, augmentation policy, and primary
evaluation rule are fixed across the main grid. The planned changes are the
feature configuration, the width multiplier, and the separately identified
secondary-backbone control.

\subsection{Feature Geometry}
Let $I_t$ and $Q_t$ denote the RMS-normalized in-phase and quadrature samples
at time index $t$, and let $\epsilon$ be the fixed feature stabilizer. The
magnitude and unit-magnitude coordinates are
\begin{equation}
\begin{aligned}
 m_t &= \sqrt{I_t^2+Q_t^2},\\
 \tilde I_t &= \frac{I_t}{\max(m_t,\epsilon)}, \qquad
 \tilde Q_t = \frac{Q_t}{\max(m_t,\epsilon)}.
\end{aligned}
\label{eq:unit-magnitude}
\end{equation}
With $\phi_t=\operatorname{atan2}(Q_t,I_t)$, define
\[
\operatorname{wrap}(x)=\operatorname{atan2}(\sin x,\cos x).
\]
The difference channels are
\begin{equation}
\begin{aligned}
 \Delta\phi_t &= \operatorname{wrap}(\phi_t-\phi_{t-1}),\\
 \Delta m_t &= m_t-m_{t-1}, \qquad
 \Delta\phi_1=\Delta m_1=0.
\end{aligned}
\label{eq:feature-differences}
\end{equation}
The stabilizers, channel order, and first-sample convention are fixed in the
feature implementation before the training grid is run. The coordinates in
Eq.~\eqref{eq:unit-magnitude} are therefore unit-magnitude coordinates, not a
second RMS normalization of $I$ and $Q$.

Four nested representations are evaluated: F2 contains $I,Q$; F3 adds $m$;
F5 further adds $\tilde I,\tilde Q$; and F7 adds $\Delta\phi,\Delta m$. The
nested ordering is a controlled study factor, not a claim that the listed
channel order is universally optimal. Deterministic feature operations are
included in the preprocessing-cost ledger and are not silently folded into
the learned-network FLOP count.

The feature ladder and its deterministic preprocessing scope are summarized in
Table~\ref{tab:feature-definitions}.
\begin{table}[!t]
\caption{Nested Feature Geometries. Preprocessing FLOPs are reported per normalized batch-one input of shape $(1,C,128)$ under the deterministic feature ledger.}
\label{tab:feature-definitions}
\centering
\footnotesize
\setlength{\tabcolsep}{1pt}
\begin{tabular}{@{}>{\raggedright\arraybackslash}p{0.11\columnwidth}>{\raggedright\arraybackslash}p{0.66\columnwidth}>{\centering\arraybackslash}p{0.19\columnwidth}@{}}
\toprule
\multicolumn{1}{>{\raggedright\arraybackslash}m{0.11\columnwidth}}{Code} &
\multicolumn{1}{>{\raggedright\arraybackslash}m{0.66\columnwidth}}{Channels} &
\multicolumn{1}{>{\raggedright\arraybackslash}m{0.19\columnwidth}}{\shortstack[c]{Preprocessing\\FLOPs}} \\
\midrule
F2 & $I,Q$ & 770 \\
F3 & $I,Q,\text{magnitude}$ & 1,410 \\
F5 & $I,Q,\text{magnitude, unit-magnitude }I/Q$ & 1,666 \\
F7 & $I,Q,\text{magnitude, unit-magnitude }I/Q,\Delta\phi,\Delta m$ & 2,429 \\
\bottomrule
\end{tabular}
\end{table}


\subsection{Backbones and Capacity}
The primary backbone, BB1, is a fixed one-dimensional depthwise-separable
convolutional network. Its stem uses a kernel of five, followed by five
depthwise-separable blocks arranged as three declared stages with base channel
widths $(24,32,48)$. A width multiplier $w$ scales all stage channels using one
common rounding rule, with $w\in\{0.50,0.75,1.00,1.50,2.00\}$. The input stem
changes only as required by the number of feature channels. At $w=1.00$, the
ledger reports 10,987 parameters for F2 and 11,587 for F7. Receiver-specific
branches, learned fusion gates, and special attention modules are excluded
from the BB1 factorization.

The resulting primary-backbone contract is summarized in
Table~\ref{tab:backbone-definition}; the cost ledger remains the authoritative
source for every feature--width cell.
\begin{table}[!t]
\caption{Primary Backbone Definition}
\label{tab:backbone-definition}
\centering
\footnotesize
\setlength{\tabcolsep}{1pt}
\begin{tabular}{@{}>{\raggedright\arraybackslash}p{0.36\columnwidth}>{\raggedright\arraybackslash}p{0.62\columnwidth}@{}}
\toprule
Item & BB1 specification \\
\midrule
Input & $(1,C,128)$, $C\in\{2,3,5,7\}$ \\
Stem & Conv1-D, kernel 5, stride 1 \\
Blocks & Five depthwise-separable residual blocks \\
Base channels & $(24,32,48)$ across three stages \\
Width set & $\{0.50,0.75,1.00,1.50,2.00\}$ \\
Rounding & Nearest multiple of four; minimum eight \\
Classifier & Global mean pooling and 11-way linear layer \\
Parameters at $w=1$ & F2: 10,987; F7: 11,587 \\
\bottomrule
\end{tabular}
\end{table}


A compact residual one-dimensional network, BB2, is used only as a directional
control at F2/F7 and the low/high width settings. BB2 is not treated as a
second full architecture benchmark; its purpose is to check whether the sign
and broad pattern of the main contrasts are specific to BB1.

\subsection{Cost Accounting and Matched Comparisons}
For each feature/backbone/width cell, the cost ledger records trainable and
total parameter counts, Core MACs, Core FLOPs, non-MAC FLOPs,
preprocessing FLOPs, Total FLOPs, input shape, profiler version, operator
coverage, and FP32 weight size. Core FLOPs are defined as twice the Core MACs.
Total FLOPs is the sum of Core FLOPs, non-MAC FLOPs, and preprocessing FLOPs.
Elementwise products, residual additions, pooling reductions, and deterministic
feature transformations not exposed by a module hook are counted explicitly.
An independent profiler is used as an arithmetic cross-check.

Matched comparisons are generated from the cost ledger before accuracy records
are read. A pair is accepted only when its declared feature or capacity
contrast is satisfied and its relative Core-FLOP difference is within the
locked tolerance. Parameters remain a secondary descriptor. No interpolated
model or artificial width is introduced when the ledger contains no valid pair.
The present manuscript reports the nominal-width confirmatory analysis and does
not use a Core-FLOP-matched pair for an accuracy claim.

\subsection{Effect Definitions}
At each width $w$, the primary feature effect is
\begin{equation}
 E_{\mathrm{feature}}(w)=A(\mathrm{F7},w)-A(\mathrm{F2},w),
\end{equation}
where $A$ denotes final-epoch accuracy under the declared split and seed unit.
F3 and F5 provide intermediate descriptive steps. For each feature
configuration $f$, the global capacity effect is
\begin{equation}
 E_{\mathrm{capacity}}(f)=A(f,2.00)-A(f,0.50).
\end{equation}
Local effects compare the adjacent 0.50-to-1.00 and 1.00-to-2.00 ranges. The
feature--capacity interaction is the change in the capacity effect across
feature configurations,
\begin{equation}
 \Delta_{\mathrm{int}} = E_{\mathrm{capacity}}(\mathrm{F7})
 - E_{\mathrm{capacity}}(\mathrm{F2}), \label{eq:interaction}
\end{equation}
with the nominal-width view reported here. The cost fields remain available for
future matched analyses, but no Core-FLOP-matched accuracy result is claimed.
These definitions are fixed before final test access.

\subsection{Statistical Analysis Plan}
The confirmatory family is defined at the split--seed unit and contains ten
paired contrasts: F7 minus F2 at 0.50, 1.00, and 2.00; the 0.50-to-2.00
capacity contrasts for F2 and F7; the 0.50-to-1.00 and 1.00-to-2.00 local
capacity contrasts for F2 and F7; and the interaction contrast in
Eq.~\eqref{eq:interaction}. Holm correction is applied across these ten
contrasts. H4 and H5 are outside the results reported in this manuscript. Paired
differences are summarized with sample size, mean difference, confidence
intervals, raw and adjusted $p$-values, and standardized effect sizes.

The primary summaries use exact sign-flip permutation $p$-values and paired
$t$ intervals. Bootstrap intervals are retained as auxiliary package outputs and are not used
for manuscript decisions. Because the study has two independently generated splits from the same corpus,
split-level inference rests on two units; seed-level intervals describe
training variability conditional on each split. A split-aware mixed model was
planned but was not fit because only two split levels were available. The
manuscript instead reports split-stratified descriptive sensitivity checks.
Conclusions are conditional on the two declared splits, consistent with the
cautions in \cite{Dietterich1998Tests}.
A zero-containing interval is described as not distinguishable under the study
precision, not as proof of equality. Additional analyses are labeled
exploratory.

\subsection{Analysis and Test Access}
Validation records are physically separated from test records. Training and
pre-unblind analysis operate on validation records only; the test-inference
command refuses to run without the study lock. Test inference is performed
only after the planned configurations, cost ledger, matching pairs, and
analysis code are frozen. Test predictions therefore cannot determine an
epoch, feature configuration, width, matched pair, or split.
